%% toc: P15 171 171.1,171.2,171.3,171.4
\chapter{Statistical Learning and Generalization}
\label{sec:p15-w02b-p15-064-prediction-risk}

\Cref{sec:p15-bounded-concentration-learning} measures a candidate by its
average loss, but only for a loss with values in \([0,1]\) attached to the
members of a class fixed in advance, and \cref{sec:p15-occam-learning}
works with the misclassification loss of a binary or finite-label problem.
Neither chapter states a prediction problem in general form, and neither
compares a learned predictor with the best predictor that exists at all.
This chapter supplies both. It states a prediction problem with arbitrary
input and output spaces and a nonnegative loss that may be unbounded, such
as squared loss for real-valued prediction; it defines risk, empirical risk
and the minimum achievable risk over all measurable predictors; it
identifies the minimizers for squared loss and for misclassification loss;
and it separates the excess risk over the best member of a class, which is
what \cref{thm:p15-conc-finite-learning} controls, from the excess risk over
the minimum achievable risk, which no sample size can reduce below the gap
between the class and that minimum.

The general set-up and the vocabulary of Bayes classifier, empirical risk,
consistency, Bayes-consistency, estimation error and approximation error
follow von Luxburg and Sch\"olkopf's presentation
\cite[Sections~2.2--2.4, pp.~2--9]{LuxburgScholkopf2008Author}, which
develops them for binary classification and mentions squared loss for
regression. The statements for an unbounded loss, the squared-loss
decomposition, the excess-risk formula for misclassification, the finite
second-moment uniform bound, the worked examples and the failure cases are
proved in full in the annex.

\emph{Inputs.} Measurable input and output spaces; a joint law of one
input--output record; a measurable loss with values in \([0,\infty)\);
an independent sample of records with that law; a class of predictors.
\emph{Conclusions.} Risk as an extended nonnegative number; empirical risk
as an unbiased estimate of the risk of a predictor fixed before sampling,
with a uniform finite-class bound under finite second moments; the minimum
achievable risk and its minimizers for squared and misclassification loss;
the decomposition of excess risk; and the exact condition under which
in-class guarantees become guarantees relative to the minimum achievable
risk.

\section{Prediction problems, loss, risk and empirical risk}
\label{sec:p15-w02b-p15-064-loss-risk}

\begin{hbdefinition}{def:p15-w02b-p15-064-prediction-problem}{Prediction problem, predictor and loss}
A \emph{prediction problem} consists of
\begin{enumerate}
\item an input space \((\mathcal X,\mathcal A_{\mathcal X})\) and an output
      space \((\mathcal Y,\mathcal A_{\mathcal Y})\), both measurable spaces;
\item a probability law \(P\) on the product space
      \(\mathcal X\times\mathcal Y\), the law of one record \((X,Y)\);
\item a measurable \emph{loss}
      \(\ell:\mathcal X\times\mathcal Y\times\mathcal Y\to[0,\infty)\),
      where \(\ell(x,y,a)\) is the cost of predicting \(a\) at input \(x\)
      when the output is \(y\).
\end{enumerate}
A \emph{predictor} is a measurable map \(f:\mathcal X\to\mathcal Y\).
The set of all predictors is \(\mathcal F_{\mathrm{all}}\); a \emph{class}
is a nonempty subset \(\mathcal F\subseteq\mathcal F_{\mathrm{all}}\).
No restriction is placed on \(P\): the output need not be a function of the
input.

Two losses recur. The \emph{squared loss}, for \(\mathcal Y=\mathbb R\)
with its Borel sets, is \(\ell(x,y,a)=(y-a)^2\); it is unbounded.
For \(\mathcal Y=\{-1,+1\}\) with all subsets measurable,
\emph{misclassification loss} (zero--one loss) is
\(\ell(x,y,a)=\mathbf1_{\{a\neq y\}}\).
\index{prediction problem}\index{loss function}\index{squared loss}\index{zero--one loss}
\end{hbdefinition}

The loss is required to be nonnegative so that its expectation is always
defined. A loss bounded below by a constant \(-c\) becomes nonnegative after
adding \(c\), which shifts every risk below by the same \(c\) and changes no
comparison between predictors. The source writes the loss as
\(\ell(X,Y,f(X))\) and adopts the convention that zero is the loss of a
perfect prediction \cite[Section~2.2, p.~4]{LuxburgScholkopf2008Author};
it identifies the predictors with the measurable maps from inputs to
outputs \cite[Section~2.2, p.~5]{LuxburgScholkopf2008Author}.

\begin{hbdefinition}{def:p15-w02b-p15-064-risk}{Risk, empirical risk and the minimum achievable risk}
In a prediction problem, the \emph{risk} of a predictor \(f\) is
\[
 R(f)=E\,\ell(X,Y,f(X))=\int\ell(x,y,f(x))\,P(dx,dy)\in[0,\infty].
\]
Let \((X_1,Y_1),\ldots,(X_n,Y_n)\), \(n\geq1\), be independent records,
each of law \(P\). The \emph{empirical risk} of \(f\) is
\[
 \widehat R_n(f)=\frac1n\sum_{i=1}^n\ell(X_i,Y_i,f(X_i)).
\]
The source writes \(R_{\mathrm{emp}}(f)\) for this quantity
\cite[Section~2.3, p.~5]{LuxburgScholkopf2008Author}.
For a class \(\mathcal F\), put
\[
 R_{\mathcal F}=\inf_{f\in\mathcal F}R(f),\qquad
 R^{*}=R_{\mathcal F_{\mathrm{all}}}=\inf_{f\in\mathcal F_{\mathrm{all}}}R(f).
\]
\(R^{*}\) is the \emph{minimum achievable risk}, or \emph{Bayes risk}, of the
problem. A predictor \(f\) with \(R(f)=R^{*}\) is a \emph{Bayes predictor};
its existence is not part of the definition. When \(R^{*}<\infty\), the
\emph{excess risk} of \(f\) is \(R(f)-R^{*}\), and its excess risk
\emph{relative to \(\mathcal F\)} is \(R(f)-R_{\mathcal F}\) whenever
\(R_{\mathcal F}<\infty\).
\index{risk}\index{empirical risk}\index{Bayes risk}\index{excess risk}
\end{hbdefinition}

A predictor \(\widehat f_n\) produced from the sample is itself random.
Its risk is then computed for a fresh record independent of the sample:
\(R(\widehat f_n)\) is the function of the sample obtained by evaluating
\(R\) at the realized predictor. When the output of a procedure is jointly
measurable in the input and the sample, this function is measurable by
Tonelli's theorem, and it is assumed below whenever a probability of an
event about \(R(\widehat f_n)\) is written.

\begin{hbdefinition}{def:p15-w02b-p15-064-consistency}{Consistency and Bayes-consistency}
Let \((X_i,Y_i)_{i\geq1}\) be independent records of law \(P\), and for each
\(n\) let \(\widehat f_n\) be the predictor a procedure computes from the
first \(n\) records, with \(R(\widehat f_n)\) measurable. Let
\(\mathcal F\) be a class with \(R_{\mathcal F}<\infty\). The procedure is
\emph{consistent with respect to \(\mathcal F\) and \(P\)} if
\(P\{R(\widehat f_n)-R_{\mathcal F}>\varepsilon\}\to0\) for every
\(\varepsilon>0\), and \emph{Bayes-consistent with respect to \(P\)} if
\(R^{*}<\infty\) and \(P\{R(\widehat f_n)-R^{*}>\varepsilon\}\to0\) for every
\(\varepsilon>0\).
\index{consistency!relative to a class}\index{Bayes-consistency}
\end{hbdefinition}

This is the source's Definition~1, parts~1 and~2
\cite[Section~2.3, p.~7]{LuxburgScholkopf2008Author}, with one change of
form. The source assumes for simplicity that a unique best member of the
class exists and compares with its risk; the infimum \(R_{\mathcal F}\) is
used here, which equals that risk whenever such a member exists. The source
also defines universal consistency, consistency for every law \(P\), and
notes that these are statements of convergence in probability.

\subsection{Empirical risk as an estimate}

\begin{hbtheorem}{thm:p15-w02b-p15-064-empirical-estimate}{proved=annex:p15-w02b-p15-064-empirical-estimate}{Unbiasedness for a fixed predictor and a finite-class bound under finite second moments}
Let \(f\) be a predictor that does not depend on the sample. Then
\(E\,\widehat R_n(f)=R(f)\) in \([0,\infty]\). If moreover
\(\sigma^2(f)=\operatorname{Var}\ell(X,Y,f(X))<\infty\), then for every
\(\varepsilon>0\)
\[
 P\{|\widehat R_n(f)-R(f)|\geq\varepsilon\}\leq\frac{\sigma^2(f)}{n\varepsilon^2}.
\]
Let \(\mathcal F=\{f_1,\ldots,f_r\}\) be a finite class fixed before sampling,
with \(\sigma^2(f_j)<\infty\) for every \(j\). Then
\[
 P\Bigl\{\max_{1\leq j\leq r}|\widehat R_n(f_j)-R(f_j)|\geq\varepsilon\Bigr\}
 \leq\frac{1}{n\varepsilon^2}\sum_{j=1}^r\sigma^2(f_j).
\]
On the complementary event, every \(\gamma\)-approximate empirical-risk
minimizer \(\widehat f\) over \(\mathcal F\), \(\gamma\geq0\), satisfies
\(R(\widehat f)<R_{\mathcal F}+2\varepsilon+\gamma\).
\end{hbtheorem}

The tolerance written \(\gamma\) here is the one written \(\eta\) in
\cref{def:p15-conc-learning}; the letter \(\eta\) is reserved in this chapter
for a conditional probability. The bound needs no range condition on the
loss, so it covers squared loss whenever \(E(Y-f_j(X))^4<\infty\) for each
\(j\), but its
dependence on the failure probability is polynomial: driving the right side
below \(\delta\) needs \(n\geq\sum_j\sigma^2(f_j)/(\delta\varepsilon^2)\),
where the bounded-loss route needs only a logarithm of \(1/\delta\). It also
needs the variances, or bounds for them, which the law of the records
determines and which are unknown in practice.

\begin{hbcorollary}{thm:p15-w02b-p15-064-bounded-rescaling}{proved=annex:p15-w02b-p15-064-bounded-rescaling}{Losses with a known finite range}
In the setting of \cref{def:p15-conc-learning}, replace the range \([0,1]\)
of the loss by \([0,B]\) for a known constant \(B>0\). For \(0<\delta<1\),
with \(b_n\) as in \cref{thm:p15-conc-finite-learning}, with probability at
least \(1-\delta\) every \(h\in\mathcal H\) satisfies
\(|\widehat R_n(h)-R(h)|\leq Bb_n\), and on the same event every
\(\eta\)-approximate minimizer \(\widehat h\) satisfies
\(R(\widehat h)\leq\min_hR(h)+2Bb_n+\eta\).
\end{hbcorollary}

Squared loss with outputs and predictions confined to an interval of length
\(L\) has range \([0,L^2]\) and is covered by this corollary with
\(B=L^2\). Without a known range, the corollary has no constant to use.

\subsection{The minimum achievable risk for two losses}

\begin{hbtheorem}{thm:p15-w02b-p15-064-squared-bayes}{proved=annex:p15-w02b-p15-064-squared-bayes}{Squared loss: the regression function and the risk decomposition}
Take \(\mathcal Y=\mathbb R\) and squared loss, and assume \(EY^2<\infty\).
Choose a measurable \(m:\mathcal X\to\mathbb R\) with \(m(X)\) a version of
\(E(Y\mid X)\), as \cref{thm:p15-conditional-factorization} permits. Then
\(E\,m(X)^2<\infty\), and for every predictor \(f\),
\[
 R(f)=E\bigl(Y-m(X)\bigr)^2+E\bigl(f(X)-m(X)\bigr)^2
\]
as an identity in \([0,\infty]\). Consequently
\(R^{*}=R(m)=E(Y-m(X))^2<\infty\), \(m\) is a Bayes predictor, and a
predictor \(f\) is a Bayes predictor exactly when \(f(X)=m(X)\) almost
surely. The excess risk of \(f\) is \(E(f(X)-m(X))^2\).
\end{hbtheorem}

\(m\) is called the regression function. The source states that the terms
bias and variance were first used for regression with squared loss
\cite[Section~2.4, p.~9]{LuxburgScholkopf2008Author}; the identity above is
the exact form of the squared-loss comparison with the best predictor and
is proved directly from the definition of conditional expectation.

\begin{hbtheorem}{thm:p15-w02b-p15-064-classification-bayes}{proved=annex:p15-w02b-p15-064-classification-bayes}{Misclassification loss and the Bayes classifier}
Take \(\mathcal Y=\{-1,+1\}\) and misclassification loss. Choose a measurable
\(\eta:\mathcal X\to[0,1]\) with \(\eta(X)\) a version of
\(P(Y=1\mid X)\), and define
\[
 f_{\mathrm{Bayes}}(x)=\begin{cases}+1,&\eta(x)\geq1/2,\\-1,&\eta(x)<1/2.\end{cases}
\]
For every predictor \(f\),
\[
 R(f)-R(f_{\mathrm{Bayes}})
 =E\bigl[\,|2\eta(X)-1|\,\mathbf1_{\{f(X)\neq f_{\mathrm{Bayes}}(X)\}}\bigr]\geq0.
\]
Hence \(f_{\mathrm{Bayes}}\) is a Bayes predictor and
\(R^{*}=E\min\{\eta(X),1-\eta(X)\}\). A predictor \(f\) is a Bayes
predictor exactly when \(f(X)=f_{\mathrm{Bayes}}(X)\) almost surely on the
event \(\{\eta(X)\neq1/2\}\). The same statements hold for labels
\(\{0,1\}\), with \(0\) in place of \(-1\) and \(\eta(X)\) a version of
\(P(Y=1\mid X)\).
\end{hbtheorem}

The function \(\eta\) and the rule \(f_{\mathrm{Bayes}}\) are the source's
equations~(1) and~(2) \cite[Section~2.2, pp.~3 and~5]{LuxburgScholkopf2008Author},
where the rule is called the Bayes classifier and described as the best
classifier that exists. The formula above gives the proof of that description
and the size of the loss incurred by disagreeing with it.

\subsection{Excess risk relative to a class and relative to the minimum}

\begin{hbtheorem}{thm:p15-w02b-p15-064-decomposition}{proved=annex:p15-w02b-p15-064-decomposition}{Estimation and approximation error, and what a class can reach}
Let \(\mathcal F\) be a class and assume \(R^{*}<\infty\) and
\(R_{\mathcal F}<\infty\). For every \(f\in\mathcal F\),
\[
 R(f)-R^{*}=\bigl(R(f)-R_{\mathcal F}\bigr)+\bigl(R_{\mathcal F}-R^{*}\bigr),
\]
and both terms are nonnegative. The second term depends only on
\(\mathcal F\), the loss and \(P\); it does not depend on the sample or on
\(n\). Consequently, for every procedure whose predictors lie in
\(\mathcal F\):
\begin{enumerate}
\item \(R(\widehat f_n)-R^{*}\geq R_{\mathcal F}-R^{*}\) for every \(n\) and
      every sample;
\item if \(R(\widehat f_n)-R_{\mathcal F}\leq\varepsilon\) on an event, then
      \(R(\widehat f_n)-R^{*}\leq\varepsilon+(R_{\mathcal F}-R^{*})\) on that
      event;
\item if \(R_{\mathcal F}>R^{*}\), the procedure is not Bayes-consistent with
      respect to \(P\); indeed
      \(P\{R(\widehat f_n)-R^{*}>\varepsilon\}=1\) for every \(n\) and every
      \(0<\varepsilon<R_{\mathcal F}-R^{*}\);
\item if \(R_{\mathcal F}=R^{*}\), then consistency with respect to
      \(\mathcal F\) and \(P\) coincides with Bayes-consistency with
      respect to \(P\).
\end{enumerate}
\end{hbtheorem}

The two terms are the source's estimation error and approximation error,
its equation~(4) \cite[Section~2.4, p.~9]{LuxburgScholkopf2008Author}.
The source attributes the first to the randomness of the sample and the
second to the choice of class; item~1 is the exact sense in which the second
is untouched by sampling.

\begin{hbcorollary}{thm:p15-w02b-p15-064-finite-class-bayes}{proved=annex:p15-w02b-p15-064-finite-class-bayes}{What the finite-class theorem says about the minimum achievable risk}
In a prediction problem, let \(\mathcal H=\{h_1,\ldots,h_r\}\) be a finite
class fixed before sampling with \(\ell(x,y,h(x))\in[0,1]\) for all
\(x,y\) and \(h\in\mathcal H\), so that \(\ell(h,(x,y))=\ell(x,y,h(x))\)
satisfies \cref{def:p15-conc-learning}. Let \(R^{*}\) be the minimum
achievable risk of the problem, and \(0<\delta<1\). On the event of
probability at least \(1-\delta\) given by
\cref{thm:p15-conc-finite-learning}, every \(\eta\)-approximate
empirical-risk minimizer \(\widehat h\) satisfies
\[
 \min_{h\in\mathcal H}R(h)-R^{*}
 \leq R(\widehat h)-R^{*}
 \leq\Bigl(\min_{h\in\mathcal H}R(h)-R^{*}\Bigr)+2b_n+\eta .
\]
Exact empirical-risk minimization over \(\mathcal H\) is consistent with
respect to \(\mathcal H\) and \(P\) for every law \(P\). It is
Bayes-consistent with respect to \(P\) if and only if some member of
\(\mathcal H\) is a Bayes predictor for \(P\).
\end{hbcorollary}

\Cref{thm:p15-conc-finite-learning} is therefore a statement about excess
risk relative to its class. Its sample-size condition guarantees an excess
over \(\min_hR(h)\) of at most \(\varepsilon\); it guarantees an excess over
\(R^{*}\) of at most \(\varepsilon\) only when the class reaches the minimum
achievable risk, and otherwise the excess over \(R^{*}\) is at least the
fixed gap \(\min_hR(h)-R^{*}\) at every sample size. The same reading applies
to \cref{thm:p15-conc-weighted-learning}, with the infimum of the penalized
risks in place of the minimum, and to the realizable statements of
\cref{sec:p15-occam-learning}, where the target concept is a member of the
class and has risk zero, so that the class reaches the minimum achievable
risk by assumption.

\subsection{Worked examples}

\begin{hbtheorem}{thm:p15-w02b-p15-064-squared-example}{proved=annex:p15-w02b-p15-064-squared-example}{A regression problem whose class misses the regression function}
Let \(X\) be uniform on \([0,1]\), let \(\xi\) be independent of \(X\) with
\(E\xi=0\) and \(E\xi^2=1\), let \(\sigma>0\), and put
\(Y=X^2+\sigma\xi\). Use squared loss. Then \(m(x)=x^2\) is a regression
function and \(R^{*}=\sigma^2\). For the class \(\mathcal F_c\) of constant
predictors \(f\equiv c\), \(c\in\mathbb R\),
\[
 R(c)=\sigma^2+\tfrac4{45}+\bigl(c-\tfrac13\bigr)^2,\qquad
 R_{\mathcal F_c}=\sigma^2+\tfrac4{45}.
\]
Empirical-risk minimization over \(\mathcal F_c\) selects the unique
minimizer \(c=\overline Y_n=n^{-1}\sum_iY_i\). Its excess risk relative to
the class is \((\overline Y_n-1/3)^2\), with expectation
\((\sigma^2+4/45)/n\), while its excess risk over \(R^{*}\) equals
\(4/45+(\overline Y_n-1/3)^2\geq4/45\) for every \(n\) and every sample.
\end{hbtheorem}

The in-class error vanishes at rate \(1/n\) in expectation, and the procedure
is consistent with respect to the constants; the distance \(4/45\) to the
minimum achievable risk is the approximation error and is the same for one
record and for a million.

\begin{hbtheorem}{thm:p15-w02b-p15-064-classification-example}{proved=annex:p15-w02b-p15-064-classification-example}{The noisy binary example of the finite-class chapter}
Take the law of \cref{thm:p15-conc-noisy-example}: \(X,Y\in\{0,1\}\),
\(P(X=0)=P(X=1)=1/2\), \(P(Y=1\mid X=0)=1/4\), \(P(Y=1\mid X=1)=3/4\), with
misclassification loss. The identity predictor \(f(x)=x\) is the Bayes
classifier and \(R^{*}=1/4\). For the class \(\mathcal H=\{0,1\}\) of the two
constant predictors, \(\min_hR(h)=1/2\). Every predictor selected from
\(\mathcal H\), by any rule, has excess risk \(0\) relative to \(\mathcal H\)
and excess risk exactly \(1/4\) over \(R^{*}\). For the four-predictor class
of \cref{thm:p15-conc-noisy-example} the two excess risks coincide.
\end{hbtheorem}

\subsection{Failure cases}

\begin{hbtheorem}{thm:p15-w02b-p15-064-failures}{proved=annex:p15-w02b-p15-064-failures}{Data-chosen predictors and unbounded risk}
\begin{enumerate}
\item \emph{Empirical risk of a data-chosen predictor.} Let \(X\) be uniform
      on \([0,1]\) and \(Y\) independent of \(X\), uniform on \(\{-1,+1\}\),
      with misclassification loss. From the sample, define
      \(\widehat f_n(x)=Y_i\) if \(x=X_i\) for some \(i\), using the smallest
      such \(i\), and \(\widehat f_n(x)=+1\) otherwise. Then almost surely
      \(\widehat R_n(\widehat f_n)=0\) and \(R(\widehat f_n)=1/2=R^{*}\).
      Thus \(E\,\widehat R_n(\widehat f_n)=0\neq1/2=E\,R(\widehat f_n)\):
      the unbiasedness of \cref{thm:p15-w02b-p15-064-empirical-estimate}
      fails for a predictor chosen with the same records.
\item \emph{Infinite minimum risk.} With squared loss, suppose \(X\) and
      \(Y\) are independent and \(EY^2=\infty\). Then \(R(f)=\infty\) for
      every predictor \(f\), so \(R^{*}=\infty\); excess risk is undefined
      and no comparison of predictors by risk is possible.
\end{enumerate}
\end{hbtheorem}

In the first case the training error is \(0\) while the risk is \(1/2\); the
source singles out an empirical risk much smaller than the true risk as the
damaging case for using the former to estimate the latter
\cite[Section~2.3, p.~6]{LuxburgScholkopf2008Author}. Here the predictor is
nonetheless a Bayes predictor, because no predictor can do better than
\(1/2\) when the output is independent of the input; the failure lies in the
estimate of its risk, not in its risk. The second case shows that the
moment condition \(EY^2<\infty\) in
\cref{thm:p15-w02b-p15-064-squared-bayes} cannot be dropped.

\begin{hbmethod}{def:p15-w02b-p15-064-method}{From a prediction problem to a statement about excess risk}
\begin{enumerate}
\item State \(\mathcal X\), \(\mathcal Y\), the loss and the sampling
      assumption; check that the loss is nonnegative and measurable.
\item Determine whether the loss has a known finite range. If so,
      \cref{thm:p15-conc-finite-learning}, rescaled by
      \cref{thm:p15-w02b-p15-064-bounded-rescaling}, or
      \cref{thm:p15-conc-weighted-learning} gives in-class guarantees with
      a logarithmic dependence on the failure probability. If not, and
      second moments are finite with known bounds,
      \cref{thm:p15-w02b-p15-064-empirical-estimate} gives a finite-class
      guarantee with polynomial dependence.
\item Report the guarantee as an excess over \(R_{\mathcal F}\). Claim an
      excess over \(R^{*}\) only with a separate argument that
      \(R_{\mathcal F}-R^{*}\) is zero or bounded, as in
      \cref{thm:p15-w02b-p15-064-decomposition}.
\item Do not use the empirical risk of a predictor selected with the same
      records as an estimate of its risk without a uniform bound covering
      the selection, or an independent validation sample as in
      \cref{thm:p15-conc-validation}.
\end{enumerate}
\end{hbmethod}

\paragraph{Relationships.}
Every in-class guarantee of \cref{sec:p15-bounded-concentration-learning}
(\cref{thm:p15-conc-finite-learning,thm:p15-conc-weighted-learning}, resting
on Hoeffding's inequalities as compared in
\cref{thm:p15-w02b-p15-063-chebyshev-comparison,thm:p15-w02b-p15-063-normal-benchmark})
and of \cref{sec:p15-occam-learning}
(\cref{thm:p15-occam-finite,thm:p15-occam-weighted,thm:p15-occam-compression})
bounds the excess over \(R_{\mathcal F}\), not over \(R^{*}\), in the sense
of \cref{thm:p15-w02b-p15-064-decomposition}; the minimum description length
principle of \cref{sec:p15-rissanen-structural-and-numerical-inference}
selects a model by total code length and is not such a risk bound.
